\begin{frame}{Background 2: chain-of-thought prompting}
\begin{block}{Extension of few-shot prompting}
Chain-of-thought adds intermediate natural-language reasoning steps $t_i$ to each demonstration.
\end{block}
\[
  p = \langle x_1 \cdot t_1 \cdot y_1 \rangle \concat \langle x_2 \cdot t_2 \cdot y_2 \rangle \concat \dots \concat \langle x_k \cdot t_k \cdot y_k \rangle
\]
\[
  p \concat x_{\text{test}} \quad \Rightarrow \quad t_{\text{test}} \cdot y_{\text{test}}
\]
\vspace{0.15cm}
\begin{columns}[T,totalwidth=\textwidth]
\begin{column}{0.48\textwidth}
\textbf{Strength:} encourages decomposition into substeps.
\end{column}
\begin{column}{0.48\textwidth}
\textbf{Weakness:} the LLM still has to execute every substep itself.
\end{column}
\end{columns}
\source{PAL paper, Section 2 and Figure 1.}
\end{frame}
